% Sección 13.3: Integrales dobles sobre regiones generales

\subsection*{Ejercicios 1--6: Evaluación de integrales iteradas sobre regiones de Tipo I y Tipo II.}
\begin{enumerate}
    \subimport{ejercicio_01/}{ejercicio.tex}
    \subimport{ejercicio_02/}{ejercicio.tex}
    \subimport{ejercicio_03/}{ejercicio.tex}
    \subimport{ejercicio_04/}{ejercicio.tex}
    \subimport{ejercicio_05/}{ejercicio.tex}
    \subimport{ejercicio_06/}{ejercicio.tex}
\end{enumerate}

\subsection*{Ejercicios 7--22: Evaluación de integrales dobles sobre regiones generales descritas matemáticamente o limitadas por curvas.}
\begin{enumerate}
    \setcounter{enumi}{6}
    \subimport{ejercicio_07/}{ejercicio.tex}
    \subimport{ejercicio_08/}{ejercicio.tex}
    \subimport{ejercicio_09/}{ejercicio.tex}
    \subimport{ejercicio_10/}{ejercicio.tex}
    \subimport{ejercicio_11/}{ejercicio.tex}
    \subimport{ejercicio_12/}{ejercicio.tex}
    \subimport{ejercicio_13/}{ejercicio.tex}
    \subimport{ejercicio_14/}{ejercicio.tex}
    \subimport{ejercicio_15/}{ejercicio.tex}
    \subimport{ejercicio_16/}{ejercicio.tex}
    \subimport{ejercicio_17/}{ejercicio.tex}
    \subimport{ejercicio_18/}{ejercicio.tex}
    \subimport{ejercicio_19/}{ejercicio.tex}
    \subimport{ejercicio_20/}{ejercicio.tex}
    \subimport{ejercicio_21/}{ejercicio.tex}
    \subimport{ejercicio_22/}{ejercicio.tex}
\end{enumerate}

\subsection*{Ejercicios 23--34: Cálculo de volúmenes de sólidos tridimensionales limitados por cilindros, paraboloides y planos.}
\begin{enumerate}
    \setcounter{enumi}{22}
    \subimport{ejercicio_23/}{ejercicio.tex}
    \subimport{ejercicio_24/}{ejercicio.tex}
    \subimport{ejercicio_25/}{ejercicio.tex}
    \subimport{ejercicio_26/}{ejercicio.tex}
    \subimport{ejercicio_27/}{ejercicio.tex}
    \subimport{ejercicio_28/}{ejercicio.tex}
    \subimport{ejercicio_29/}{ejercicio.tex}
    \subimport{ejercicio_30/}{ejercicio.tex}
    \subimport{ejercicio_31/}{ejercicio.tex}
    \subimport{ejercicio_32/}{ejercicio.tex}
    \subimport{ejercicio_33/}{ejercicio.tex}
    \subimport{ejercicio_34/}{ejercicio.tex}
\end{enumerate}

\subsection*{Ejercicios 35--40: Bosquejo de regiones e inversión del orden de integración.}
\begin{enumerate}
    \setcounter{enumi}{34}
    \subimport{ejercicio_35/}{ejercicio.tex}
    \subimport{ejercicio_36/}{ejercicio.tex}
    \subimport{ejercicio_37/}{ejercicio.tex}
    \subimport{ejercicio_38/}{ejercicio.tex}
    \subimport{ejercicio_39/}{ejercicio.tex}
    \subimport{ejercicio_40/}{ejercicio.tex}
\end{enumerate}

\subsection*{Ejercicios 41--46: Cálculo de integrales iteradas mediante inversión del orden de integración.}
\begin{enumerate}
    \setcounter{enumi}{40}
    \subimport{ejercicio_41/}{ejercicio.tex}
    \subimport{ejercicio_42/}{ejercicio.tex}
    \subimport{ejercicio_43/}{ejercicio.tex}
    \subimport{ejercicio_44/}{ejercicio.tex}
    \subimport{ejercicio_45/}{ejercicio.tex}
    \subimport{ejercicio_46/}{ejercicio.tex}
\end{enumerate}

\subsection*{Ejercicios 47--48: Descomposición de regiones e integración con gráficos vectoriales.}
\begin{enumerate}
    \setcounter{enumi}{46}
    \subimport{ejercicio_47/}{ejercicio.tex}
    \subimport{ejercicio_48/}{ejercicio.tex}
\end{enumerate}

\subsection*{Ejercicios 49--53: Estimación de integrales mediante propiedades, sumas de integrales y simetría.}
\begin{enumerate}
    \setcounter{enumi}{48}
    \subimport{ejercicio_49/}{ejercicio.tex}
    \subimport{ejercicio_50/}{ejercicio.tex}
    \subimport{ejercicio_51/}{ejercicio.tex}
    \subimport{ejercicio_52/}{ejercicio.tex}
    \subimport{ejercicio_53/}{ejercicio.tex}
\end{enumerate}
